\section*{Capítulo \ref{ch:g12:plant-rooted-life} --- A vida vegetal: prosperar enraizado}

\begin{solution}{exo:g12:plant-rooted-life:1}
O sistema radicular, cujos \omterm{prop:g12:plant-rooted-life:surfaces}{pelos absorventes} apresentam uma vasta
superfície ao solo, e o sistema caulinar, cujas folhas apresentam uma
vasta superfície à luz e ao ar.
\end{solution}

\begin{solution}{exo:g12:plant-rooted-life:2}
Um poro na pele da folha delimitado por duas células-guarda. O gás
carbônico entra, o vapor de água (e o oxigênio) sai. Aberto na luz,
fechado à noite e na seca.
\end{solution}

\begin{solution}{exo:g12:plant-rooted-life:3}
Seiva bruta: água e íons minerais, no \omterm{prop:g12:plant-rooted-life:saps}{xilema}, das raízes às folhas,
puxada pela transpiração. Seiva elaborada: água e açúcares, no \omterm{prop:g12:plant-rooted-life:saps}{floema},
das folhas aos órgãos que usam ou armazenam açúcar, nas duas direções,
empurrada pela pressão das \omterm{def:g10:cells-common-unit:cell}{células} carregadas de açúcar.
\end{solution}

\begin{solution}{exo:g12:plant-rooted-life:4}
Estruturais: cutícula e casca, espinhos (também pelos urticantes,
sílica, \omterm{def:g10:cells-common-unit:organelle}{paredes celulares}). Químicas: taninos, alcaloides (também
toxinas, resinas, látex).
\end{solution}

\begin{solution}{exo:g12:plant-rooted-life:5}
Uma resposta de crescimento dirigida a um estímulo (luz, gravidade,
água). A auxina, feita na ponta em crescimento do ramo.
\end{solution}

\begin{solution}{exo:g12:plant-rooted-life:6}
$\qty{40}{cm^2} = \qty{4000}{mm^2}$: $4000 \times 200 = \num{8e5}$
\omterm{def:g12:plant-rooted-life:stomata}{estômatos} por folha; $\num{1.6e11}$ na árvore.
\end{solution}

\begin{solution}{exo:g12:plant-rooted-life:7}
Por volta do meio-dia e do começo da tarde. No calor da tarde, as
células-guarda fecham os \omterm{def:g12:plant-rooted-life:stomata}{estômatos} em parte para limitar a perda de
água, e a transpiração acompanha a abertura, não a luz.
\end{solution}

\begin{solution}{exo:g12:plant-rooted-life:8}
As folhas ficam verdes e abastecidas de água (o \omterm{prop:g12:plant-rooted-life:saps}{xilema}, dentro da
madeira, está intacto). As raízes, cortadas do açúcar das folhas,
esgotam as reservas delas ao longo de meses e morrem, e a árvore com
elas. Logo acima do anel, a casca incha com o açúcar que não consegue
mais passar.
\end{solution}

\begin{solution}{exo:g12:plant-rooted-life:9}
Intacta: a resposta existe. Ponta coberta: é na ponta que a luz é
percebida. Base coberta: a base não precisa ser iluminada para se
curvar --- ela responde a um sinal, não à luz. Ponta removida: a ponta
é necessária; o sinal vem dela. Juntos, eles separam a percepção
(ponta) da resposta (base) e mostram que algo viaja entre as duas.
\end{solution}

\begin{solution}{exo:g12:plant-rooted-life:10}
Do sol: a evaporação da água na superfície da folha, impulsionada pelo
calor solar e pelo ar seco, puxa a coluna contínua de água \omterm{prop:g12:plant-rooted-life:saps}{xilema}
acima. A planta fornece os tubos e os \omterm{def:g12:plant-rooted-life:stomata}{estômatos} abertos, não a energia.
\end{solution}

\begin{solution}{exo:g12:plant-rooted-life:11}
Flor de vento: pequena, sem graça, sem aroma, sem néctar, quantidades
enormes de pólen leve, estigmas plumosos --- o vento é cego e de graça,
e assim a planta investe na quantidade. Flor de inseto: grande,
colorida, perfumada, com néctar, pólen pegajoso em quantidades modestas
--- o transportador precisa ser atraído e pago, e entrega o pólen com precisão.
\end{solution}

\begin{solution}{exo:g12:plant-rooted-life:12}
Ela ganha água, mantendo as \omterm{def:g10:cells-common-unit:cell}{células} dela túrgidas e vivas; ela perde a
entrada de gás carbônico, portanto a \omterm{def:g10:cell-metabolism:autotrophy}{fotossíntese}, e o resfriamento da
transpiração. \omterm{def:g12:plant-rooted-life:stomata}{Estômatos} fechados significam nenhum açúcar produzido: a
planta vive das reservas dela, e uma seca longa a mata de fome antes de
a secar.
\end{solution}

\begin{solution}{exo:g12:plant-rooted-life:13}
Na folha, uma defesa: uma toxina contra os herbívoros, embutida no
tecido. No néctar, um recrutamento: uma dose baixa que faz o
polinizador lembrar e voltar, paga junto com o açúcar. Uma substância,
duas necessidades, dosada conforme o lugar.
\end{solution}

\begin{solution}{exo:g12:plant-rooted-life:14}
As flores de tubo um pouco mais longo são polinizadas apenas pelas
mariposas que chegam ao fundo, e assim o pólen delas vai para flores do
mesmo tipo; as mariposas de língua um pouco mais longa conseguem néctar
que as outras não conseguem; cada pequena vantagem é selecionada, e o
tubo e a língua se alongam juntos ao longo de muitas gerações. Se a
mariposa desaparecer, a flor não produz semente; se a flor desaparecer,
a mariposa perde o único alimento dela.
\end{solution}

\begin{solution}{exo:g12:plant-rooted-life:15}
Uma planta percebe a luz com a ponta dos ramos dela e a gravidade com
as raízes dela; ela responde crescendo em direção a uma e ao longo da
outra; ela se defende com casca e espinhos, e fabrica toxinas e
bloqueadores da digestão poucas horas depois de um ataque; e ela
recruta insetos para carregar o pólen dela, aves para carregar as
sementes dela e vespas para matar as lagartas dela, pagando em néctar,
fruto e aroma. Ela age sem se mover.
\end{solution}

\begin{solution}{pb:g12:plant-rooted-life:1}
\textbf{1.} $\num{250000} \times \qty{25}{cm^2} = \num{6.25e6}\,
\unit{cm^2} = \qty{625}{m^2}$.

\textbf{2.} $\qty{625}{m^2} = \num{6.25e8}\,\unit{mm^2}$; $\times 250
\approx \num{1.6e11}$ \omterm{def:g12:plant-rooted-life:stomata}{estômatos}.

\textbf{3.} $625/150 \approx 4$ metros quadrados de folha por metro
quadrado de chão.

\textbf{4.} Cerca de \qty{1000}{m^2} de \omterm{prop:g12:plant-rooted-life:surfaces}{pelos absorventes} contra
\qty{625}{m^2} de folhas: comparáveis. As duas são superfícies de troca
dimensionadas para os mesmos fluxos --- a água que as folhas perdem
precisa ser absorvida pelas raízes.

\textbf{5.} A luz é absorvida em uma fração de milímetro de tecido, e o
gás carbônico se difunde apenas por distâncias curtas: um órgão espesso
deixaria a maioria das \omterm{def:g10:cells-common-unit:cell}{células} dele no escuro e sem alimento. Folhas
finas e empilhadas põem cada \omterm{def:g10:cells-common-unit:cell}{célula} perto da luz e do ar.

\textbf{6.} $625 \times 0.6 = \qty{375}{L}$.

\textbf{7.} $625 \times 8 = \qty{5}{kg}$ de gás carbônico, o que dá
cerca de $5/1.47 \approx \qty{3.4}{kg}$ de açúcar (ou, mais simples: o
açúcar feito a partir de \qty{5}{kg} de $\mathrm{CO_2}$ é cerca de
\qty{3.4}{kg}); água usada $3.4/1.6 \approx \qty{2.1}{kg}$: cerca de
0.6\% dos \qty{375}{L} transpirados.

\textbf{8.} O gás carbônico entra pelos mesmos \omterm{def:g12:plant-rooted-life:stomata}{estômatos} abertos que
deixam a água sair; fechá-los interrompe a \omterm{def:g10:cell-metabolism:autotrophy}{fotossíntese}. A água é o
preço do açúcar.

\textbf{9.} O sol, ao evaporar a água nas folhas. Levantar
\qty{375}{kg} por \qty{20}{m} são cerca de \qty{75}{kJ} de trabalho ---
pouco em si, mas uma bomba biológica exigiria tecidos, ATP e
manutenção que a árvore não precisa construir.

\textbf{10.} Ela fechou em parte os \omterm{def:g12:plant-rooted-life:stomata}{estômatos} contra o calor seco; a
entrada de gás carbônico e a produção de açúcar caem numa fração
parecida --- uns \qty{2}{kg} de açúcar perdidos ao longo da tarde.

\textbf{11.} Cerca de \qty{3.4}{kg} produzidos; \qty{2}{kg} (60\%) para
as raízes, o tronco e os ramos, \qty{0.5}{kg} (15\%) para as bolotas.

\textbf{12.} O \omterm{prop:g12:plant-rooted-life:saps}{floema}, das folhas (fonte) para as raízes, o tronco, os
ramos e as bolotas (drenos).

\textbf{13.} Armazenado como amido no tronco e nas raízes; na primavera
seguinte ele é convertido de volta em açúcar para construir as folhas
novas antes que elas consigam se alimentar sozinhas.

\textbf{14.} O açúcar continua chegando aos ramos e às bolotas acima do
anel; as raízes e o tronco de baixo não recebem nada: a parcela de
\qty{2}{kg} deles cai ao que estiver armazenado, e eles passam fome ao
longo da temporada.

\textbf{15.} As folhas restantes continuam produzindo açúcar, o amido
armazenado do tronco e das raízes cobre a falta, e a árvore consegue
lançar folhas novas a partir dos muitos pontos de crescimento dela; as
bolotas são um dreno que a árvore ainda consegue encher se o açúcar do
verão bastar.

\textbf{16.} Flores pequenas e esverdeadas, sem pétalas, aroma nem
néctar, penduradas em amentilhos que soltam nuvens de pólen leve;
estigmas grandes e plumosos para capturá-lo.

\textbf{17.} As bolotas largadas sob o genitor crescem, se é que
crescem, na sombra dele e competem com ele; um gaio as leva centenas de
metros e as enterra na profundidade certa. Perder a maioria delas é o
preço de colocar algumas bem.

\textbf{18.} $3000 \times 0.05 = 150$ bolotas esquecidas; $150/50 = 3$
mudas.

\textbf{19.} Três mudas por ano ao longo de algumas décadas de boas
safras dão talvez uma centena de mudas na vida da árvore, das quais uma
precisa chegar à idade adulta. Os números são grandes porque quase toda
semente e toda muda se perdem.

\textbf{20.} \qty{625}{m^2} de folha e uns $\num{1.6e11}$ \omterm{def:g12:plant-rooted-life:stomata}{estômatos};
cerca de \qty{375}{L} levantados num dia por \qty{3.4}{kg} de açúcar; o
vento para o pólen dele e o gaio para as bolotas dele.
\end{solution}
